\clearpage
\begin{proofbox}
\textbf{路线一：匹配导数，最快确定系数。}
若 $P_n(x)=a_0+a_1x+\cdots+a_nx^n$，则
$P_n^{(k)}(0)=k!a_k$。要求它与 $f$ 在 $0$ 处的各阶导数相同，只能取
\[
 a_k=\frac{f^{(k)}(0)}{k!}\qquad(0\le k\le n).
\]
这是确定系数的办法；还需证明余项小，不能把“匹配”本身当成误差证明。

\textbf{一个低门槛的误差工具：只用导数比较。}
设整数 $m\ge1$、$r>0$，$g\in C^m([-r,r])$，且
$g(0)=g'(0)=\cdots=g^{(m-1)}(0)=0$、$|g^{(m)}|\le M$，则
\[
 |g(x)|\le\frac{M|x|^m}{m!}\qquad(|x|\le r).
\]
证明：先取 $0\le t\le r$。由
$-M\le g^{(m)}(t)\le M$ 及单调性比较，得
$-Mt\le g^{(m-1)}(t)\le Mt$。
若 $m=1$，已得所需结论；若 $m\ge2$，再将 $g^{(m-2)}$ 与 $\pm Mt^2/2$ 比较其导数，得到
$|g^{(m-2)}(t)|\le Mt^2/2$。
依次比较到 $g$，每次幂次加一并除以新幂次，得到所需阶乘。
负半轴对 $h(t)=g(-t)$ 重复同一论证。因此左右两侧都覆盖。

\textbf{指数：三次导数直接控制二次近似。}
取 $g=e^x-E_2(x)$，其 $0,1,2$ 阶导数在 $0$ 处都为 $0$，且
$g^{(3)}(x)=e^x$。在 $[-r,r]$ 上 $|g^{(3)}|\le e^r$，故
\[
 |e^x-E_2(x)|\le e^r|x|^3/6.
\]
固定 $r$ 后，这就严格给出 $O(x^3)$。

\textbf{正弦、余弦：多检查一个零系数，避免余项跳阶。}
弧度制下导数循环为
\[
\begin{array}{c|rrrrrrrr}
 k&0&1&2&3&4&5&6&7\\ \hline
 (\sin x)^{(k)}|_{x=0}&0&1&0&-1&0&1&0&-1\\
 (\cos x)^{(k)}|_{x=0}&1&0&-1&0&1&0&-1&0
\end{array}
\]
对 $g=\sin x-S_5(x)$，零点处直到\textbf{第六阶}导数都为零，第七阶导数为 $-\cos x$，绝对值不超过 $1$。
对 $g=\cos x-C_4(x)$，零点处直到\textbf{第五阶}导数都为零，第六阶导数为 $-\cos x$，绝对值也不超过 $1$。
应用上述工具即得 $|x|^7/7!$ 与 $|x|^6/6!$ 两个界。
这说明余项分别为 $O(x^7)$、$O(x^6)$，而非仅凭“奇偶性”省略余项论证。

\tcbbreak
\textbf{路线二：单位圆先解释最基本的三角近似。}
\begin{center}
\begin{tikzpicture}[font=\small]
\draw[black!30,rounded corners=2pt] (0,0) rectangle (14.4,8.2);
\node[anchor=west,font=\small\bfseries] at (0.5,7.72) {弧度的几何含义：单位圆上，弧长就是角的数值};

\pgfmathsetmacro{\Ffourangle}{0.65*180/pi}
\begin{scope}[shift={(3.7,3.6)},scale=2.45]
\fill[purple!8] (0,0)--(1,0)--(1,{tan(\Ffourangle)})--cycle;
\fill[orange!15] (0,0)--(1,0) arc[start angle=0,end angle=\Ffourangle,radius=1]--cycle;
\fill[blue!12] (0,0)--(1,0)--({cos(\Ffourangle)},{sin(\Ffourangle)})--cycle;
\draw[blue!60!black,thick] (1,0)--({cos(\Ffourangle)},{sin(\Ffourangle)});
\draw[black!50] (0,0) circle (1);
\draw[-Stealth] (-1.15,0)--(1.3,0);
\draw[-Stealth] (0,-1.1)--(0,1.2);
\draw[thick] (0,0)--(1,{tan(\Ffourangle)});
\draw[blue,very thick] (0,0)--({cos(\Ffourangle)},{sin(\Ffourangle)});
\draw[blue,very thick] ({cos(\Ffourangle)},0)--({cos(\Ffourangle)},{sin(\Ffourangle)});
\draw[orange!85!black,very thick] (1,0) arc[start angle=0,end angle=\Ffourangle,radius=1];
\draw[purple,very thick] (1,0)--(1,{tan(\Ffourangle)});
\node[below left] at (0,0) {$O$};
\node[below] at (1,0) {$A$};
\node[above left] at ({cos(\Ffourangle)},{sin(\Ffourangle)}) {$P$};
\node[above right] at (1,{tan(\Ffourangle)}) {$T$};
\node[below] at ({cos(\Ffourangle)},0) {$Q$};
\draw (0.3,0) arc[start angle=0,end angle=\Ffourangle,radius=.3];
\node at (0.45,0.16) {$x$};
\end{scope}
\node[anchor=west] at (8,6.2) {$0<x<\pi/2,\quad OA=1$};
\node[anchor=west,text=blue] at (8,5.4) {$PQ=\sin x$};
\node[anchor=west,text=orange!85!black] at (8,4.6) {$\widehat{AP}\text{ 的弧长}=x$};
\node[anchor=west,text=purple] at (8,3.8) {$AT=\tan x$};
\node[anchor=west] at (8,2.6) {$\sin x<x<\tan x$};
\node[anchor=west] at (8,1.7) {图示取 $x=0.65$ 弧度};

\end{tikzpicture}
\end{center}

对 $0<x<\pi/2$，三角形 $OAP$ 在扇形 $OAP$ 内，扇形又在三角形 $OAT$ 内。
比较三者面积：
\[
 \tfrac12\sin x<\tfrac12 x<\tfrac12\tan x.
\]
于是 $\cos x<\sin x/x<1$，夹逼得 $\sin x/x\to1$；负半轴由奇性得到相同结论。
再由 $\tan x/x=(\sin x/x)/\cos x$ 得 $\tan x/x\to1$。
几何路线解释一次近似和弧度的必要性，高阶系数仍由前述导数或下述代数运算确定。

\textbf{路线三：对数用有限恒等式，不用求四次导数。}
对 $t>-1$，多项式除法给出精确等式
\[
 \frac1{1+t}=1-t+t^2-\frac{t^3}{1+t}.
\]
令 $D(x)=L_3(x)-\ln(1+x)$，则 $D(0)=0$ 且
\[
 D'(x)=\frac{x^3}{1+x}.
\]
当 $|x|\le r<1$，有 $|D'(x)|\le |x|^3/(1-r)$。
对正半轴逐点比较 $D$ 与 $\pm x^4/[4(1-r)]$ 的导数；
对负半轴令 $x=-t$ 作同样比较，得到
\[
 |D(x)|\le\frac{|x|^4}{4(1-r)}.
\]
又因 $D'$ 在 $(-1,0)$ 为负、在 $(0,+\infty)$ 为正，$D$ 在 $0$ 处取最小值 $0$。
所以 $\ln(1+x)\le L_3(x)$ 对 $x>-1$ 都成立，等号仅在 $x=0$。
这既证明三次系数，也证明 $O(x^4)$ 及误差方向。

\begin{center}
\begin{tikzpicture}[font=\small]
\draw[black!30,rounded corners=2pt] (0,0) rectangle (14.4,8.2);
\node[anchor=west,font=\small\bfseries] at (0.5,7.72) {对数的误差：上方多项式导数减去下方真实导数};
\begin{axis}[width=11.2cm,height=6.1cm,at={(2.05cm,1.25cm)},anchor=south west,
axis lines=left,xlabel={$x$},ylabel={$y$},tick label style={font=\small},label style={font=\small},
legend style={font=\small,draw=none,fill=white,at={(0.03,0.97)},anchor=north west},
legend cell align=left,unbounded coords=jump,xmin=0,xmax=0.8,ymin=0.45,ymax=1.3]
\addplot[draw=none,fill=orange!20,forget plot] coordinates {(0,1) (0.01,0.9901) (0.02,0.9804) (0.03,0.9709) (0.04,0.9616) (0.05,0.9525) (0.06,0.9436) (0.07,0.9349) (0.08,0.9264) (0.09,0.9181) (0.1,0.91) (0.11,0.9021) (0.12,0.8944) (0.13,0.8869) (0.14,0.8796) (0.15,0.8725) (0.16,0.8656) (0.17,0.8589) (0.18,0.8524) (0.19,0.8461) (0.2,0.84) (0.21,0.8341) (0.22,0.8284) (0.23,0.8229) (0.24,0.8176) (0.25,0.8125) (0.26,0.8076) (0.27,0.8029) (0.28,0.7984) (0.29,0.7941) (0.3,0.79) (0.31,0.7861) (0.32,0.7824) (0.33,0.7789) (0.34,0.7756) (0.35,0.7725) (0.36,0.7696) (0.37,0.7669) (0.38,0.7644) (0.39,0.7621) (0.4,0.76) (0.41,0.7581) (0.42,0.7564) (0.43,0.7549) (0.44,0.7536) (0.45,0.7525) (0.46,0.7516) (0.47,0.7509) (0.48,0.7504) (0.49,0.7501) (0.5,0.75) (0.51,0.7501) (0.52,0.7504) (0.53,0.7509) (0.54,0.7516) (0.55,0.7525) (0.56,0.7536) (0.57,0.7549) (0.58,0.7564) (0.59,0.7581) (0.6,0.76) (0.61,0.7621) (0.62,0.7644) (0.63,0.7669) (0.64,0.7696) (0.65,0.7725) (0.66,0.7756) (0.67,0.7789) (0.68,0.7824) (0.69,0.7861) (0.7,0.79) (0.71,0.7941) (0.72,0.7984) (0.73,0.8029) (0.74,0.8076) (0.75,0.8125) (0.76,0.8176) (0.77,0.8229) (0.78,0.8284) (0.79,0.8341) (0.8,0.84) (0.8,0.5555555555555555555555555556) (0.79,0.5586592178770949720670391061) (0.78,0.5617977528089887640449438202) (0.77,0.5649717514124293785310734463) (0.76,0.5681818181818181818181818182) (0.75,0.5714285714285714285714285714) (0.74,0.5747126436781609195402298851) (0.73,0.5780346820809248554913294798) (0.72,0.5813953488372093023255813953) (0.71,0.5847953216374269005847953216) (0.7,0.5882352941176470588235294118) (0.69,0.5917159763313609467455621302) (0.68,0.5952380952380952380952380952) (0.67,0.5988023952095808383233532934) (0.66,0.6024096385542168674698795181) (0.65,0.6060606060606060606060606061) (0.64,0.6097560975609756097560975610) (0.63,0.6134969325153374233128834356) (0.62,0.6172839506172839506172839506) (0.61,0.6211180124223602484472049689) (0.6,0.625) (0.59,0.6289308176100628930817610063) (0.58,0.6329113924050632911392405063) (0.57,0.6369426751592356687898089172) (0.56,0.6410256410256410256410256410) (0.55,0.6451612903225806451612903226) (0.54,0.6493506493506493506493506494) (0.53,0.6535947712418300653594771242) (0.52,0.6578947368421052631578947368) (0.51,0.6622516556291390728476821192) (0.5,0.6666666666666666666666666667) (0.49,0.6711409395973154362416107383) (0.48,0.6756756756756756756756756757) (0.47,0.6802721088435374149659863946) (0.46,0.6849315068493150684931506849) (0.45,0.6896551724137931034482758621) (0.44,0.6944444444444444444444444444) (0.43,0.6993006993006993006993006993) (0.42,0.7042253521126760563380281690) (0.41,0.7092198581560283687943262411) (0.4,0.7142857142857142857142857143) (0.39,0.7194244604316546762589928058) (0.38,0.7246376811594202898550724638) (0.37,0.7299270072992700729927007299) (0.36,0.7352941176470588235294117647) (0.35,0.7407407407407407407407407407) (0.34,0.7462686567164179104477611940) (0.33,0.7518796992481203007518796992) (0.32,0.7575757575757575757575757576) (0.31,0.7633587786259541984732824427) (0.3,0.7692307692307692307692307692) (0.29,0.7751937984496124031007751938) (0.28,0.78125) (0.27,0.7874015748031496062992125984) (0.26,0.7936507936507936507936507937) (0.25,0.8) (0.24,0.8064516129032258064516129032) (0.23,0.8130081300813008130081300813) (0.22,0.8196721311475409836065573770) (0.21,0.8264462809917355371900826446) (0.2,0.8333333333333333333333333333) (0.19,0.8403361344537815126050420168) (0.18,0.8474576271186440677966101695) (0.17,0.8547008547008547008547008547) (0.16,0.8620689655172413793103448276) (0.15,0.8695652173913043478260869565) (0.14,0.8771929824561403508771929825) (0.13,0.8849557522123893805309734513) (0.12,0.8928571428571428571428571429) (0.11,0.9009009009009009009009009009) (0.1,0.9090909090909090909090909091) (0.09,0.9174311926605504587155963303) (0.08,0.9259259259259259259259259259) (0.07,0.9345794392523364485981308411) (0.06,0.9433962264150943396226415094) (0.05,0.9523809523809523809523809524) (0.04,0.9615384615384615384615384615) (0.03,0.9708737864077669902912621359) (0.02,0.9803921568627450980392156863) (0.01,0.9900990099009900990099009901) (0,1)} \closedcycle;
\addplot[blue,very thick,domain=0:0.8,samples=100] {1/(1+x)};\addlegendentry{$1/(1+x)$}
\addplot[orange!85!black,thick,dashed,domain=0:0.8] {1-x+x^2};\addlegendentry{$1-x+x^2$}
\node[anchor=west] at (axis cs:0.26,1.13) {橙色面积 $=L_3(0.8)-\ln1.8$};
\end{axis}

\end{tikzpicture}
\end{center}

\textbf{会积分时的另一种读法：}
\[
 L_3(x)-\ln(1+x)=\int_0^x\frac{t^3}{1+t}\,\mathrm dt.
\]
图中取 $x=0.8$，两条导数曲线之间的面积正好是函数值之差；
负 $x$ 时须同时考虑积分方向，不能把“带符号积分”直接当成正向面积。

\textbf{路线四：正切用正弦除以余弦，少求高阶导数。}
由已证公式及降阶，
\[
 \sin x=x-\frac{x^3}{6}+O(x^5),\qquad
 \cos x=1-\frac{x^2}{2}+O(x^4).
\]
令 $u=\cos x-1=-x^2/2+O(x^4)=O(x^2)$。
因为 $u\to0$，$1+u$ 在零点附近远离 $0$，利用精确恒等式
\[
 \frac1{1+u}=1-u+\frac{u^2}{1+u}
             =1+\frac{x^2}{2}+O(x^4).
\]
相乘并保留到三次：
\[
\begin{aligned}
 \tan x&=\left(x-\frac{x^3}{6}+O(x^5)\right)
          \left(1+\frac{x^2}{2}+O(x^4)\right)\\
 &=x+\left(\frac12-\frac16\right)x^3+O(x^5)
  =x+\frac{x^3}{3}+O(x^5).
\end{aligned}
\]
其中 $x\,O(x^4)=O(x^5)$，$x^3x^2=O(x^5)$；这些是五阶余项的具体来源。
五条展开和已列误差界至此全部得到证明，不需要先背无限级数。
\end{proofbox}
